\section{José Ángel Medina Alcazar}
El estudio de la robótica desde una perspectiva teórica permite comprender las bases que hacen posible el diseño y la operación de sistemas autónomos. El análisis de la cinemática directa e inversa proporciona herramientas para definir y controlar el movimiento del robot, mientras que la dinámica permite anticipar y compensar las fuerzas internas y externas que afectan su comportamiento. Por otro lado, los sistemas de control aseguran precisión y estabilidad en la ejecución de tareas complejas. Finalmente, plataformas como ROS han demostrado ser fundamentales para integrar hardware y software de forma modular, lo cual facilita el desarrollo, simulación y prueba de algoritmos robóticos tanto en entornos reales como virtuales. Estos fundamentos son esenciales no solo para implementar soluciones funcionales, sino también para avanzar en la investigación y mejora de la robótica.